% !TEX program = xelatex
% Compile: xelatex physics_ch10_urdu.tex  (run twice)
\documentclass[12pt]{article}

% --- normal packages BEFORE polyglossia (polyglossia loads bidi last) ---
\usepackage{fontspec}
\usepackage{amsmath}
\usepackage{amssymb}
\usepackage{geometry}
\geometry{a4paper, margin=2.2cm}
\usepackage{array}
\usepackage{booktabs}
\usepackage{xcolor}
\usepackage{tikz}
\usetikzlibrary{arrows.meta,decorations.markings,angles,quotes,patterns,calc}

\definecolor{seccolor}{RGB}{30,90,140}
\definecolor{rulecolor}{RGB}{205,150,35}
\definecolor{accent}{RGB}{180,60,40}
\definecolor{accent2}{RGB}{40,120,90}
\definecolor{lightfill}{RGB}{232,240,247}

\setmainfont{Latin Modern Roman}

\usepackage{polyglossia}
\setdefaultlanguage{urdu}
\setotherlanguage{english}
\newfontfamily\urdufont[Script=Arabic,Language=Urdu]{Noto Nastaliq Urdu}
\newfontfamily\englishfont{Latin Modern Roman}

\linespread{2.0}
\setlength{\parskip}{6pt}
\setlength{\parindent}{0pt}

\renewcommand{\emph}[1]{\textbf{#1}}

\newcommand{\bigheading}[1]{%
  \par\vspace{10pt}{\color{seccolor}\LARGE\bfseries #1}\par
  {\color{rulecolor}\rule{\linewidth}{1.2pt}}\par\vspace{4pt}}
\newcommand{\heading}[1]{%
  \par\vspace{6pt}{\color{seccolor}\Large\bfseries #1}\par\vspace{2pt}}
\newcommand{\subheading}[1]{%
  \par\vspace{4pt}{\large\bfseries #1}\par}

% inline LTR math/number helper
\newcommand{\m}[1]{\textenglish{$#1$}}

% figure caption (Urdu, RTL normal flow)
\newcommand{\figcap}[1]{\par\vspace{1pt}{\small\color{seccolor}#1}\par\vspace{4pt}}

\begin{document}

% ===================== TITLE =====================
\begin{center}
{\color{seccolor}\Huge\bfseries طبیعیات}\\[10pt]
{\LARGE باب ۱۰ — گردش}\\[6pt]
{\large (اردو میں تدریسی مواد، خاکوں کے ساتھ)}
\end{center}
{\color{rulecolor}\rule{\linewidth}{1.4pt}}

\vspace{4pt}

اب تک ہم نے زیادہ تر \textbf{انتقالی حرکت} (\textenglish{translation}) کا مطالعہ کیا، جس میں کوئی جسم کسی سیدھی یا خمدار لکیر کے ساتھ آگے بڑھتا ہے۔ اِس باب میں ہم \textbf{گردشی حرکت} (\textenglish{rotation}) کی طرف آتے ہیں، جس میں کوئی جسم کسی محور کے گرد گھومتا ہے۔ دلچسپ بات یہ ہے کہ گردش کے متغیرات یک بُعدی حرکت کے متغیرات سے بہت مشابہ ہیں — اگر کسی سوال میں ”زاویائی“ کا لفظ ہٹا دیا جائے تو اکثر وہ باب ۲ کا عام حرکت کا سوال بن جاتا ہے۔

% ===================== 10-1 =====================
\bigheading{۱۰-۱ \ \ گردشی متغیرات}

\heading{صلب جسم اور محورِ گردش}

\textbf{صلب جسم} (\textenglish{rigid body}) وہ جسم ہے جو اپنے تمام حصوں سمیت، اپنی شکل بدلے بغیر، یکجا ہو کر گھوم سکے۔ \textbf{مقررہ محور} (\textenglish{fixed axis}) سے مراد ایسا محور ہے جو خود حرکت نہ کرے۔ خالص گردش میں جسم کا ہر نقطہ ایک دائرے میں گھومتا ہے جس کا مرکز محورِ گردش پر ہوتا ہے، اور ہر نقطہ کسی مقررہ وقفے میں یکساں زاویہ طے کرتا ہے۔

\heading{زاویائی مقام}

جسم میں جڑی ایک \textbf{خطِ حوالہ} (\textenglish{reference line}) فرض کریں جو محور کے عمود پر ہو اور جسم کے ساتھ گھومے۔ اِس خط کا کسی مقررہ سمت (مثلاً مثبت \textenglish{$x$} محور) سے بننے والا زاویہ \textbf{زاویائی مقام} \m{\theta} کہلاتا ہے۔ اگر دائرے کا نصف قطر \m{r} اور قوس کی لمبائی \m{s} ہو تو ریڈین میں:
\[ \theta = \frac{s}{r} \qquad (\text{radian measure}). \]

\begin{center}
\begin{LTR}
\begin{tikzpicture}[scale=1.0]
  % axes
  \draw[->,gray!60] (-0.4,0)--(4.6,0) node[right]{\small $x$};
  \draw[->,gray!60] (0,-0.4)--(0,4.2) node[above]{\small $y$};
  % circle
  \draw[thin,gray!50] (0,0) circle (3.2);
  % radius / reference line at 50 deg
  \draw[very thick,accent] (0,0)--(50:3.2) node[midway,above,sloped]{\small $r$};
  \node[accent] at (50:3.55) {\small reference line};
  % arc s
  \draw[very thick,accent2] (3.2,0) arc (0:50:3.2);
  \node[accent2] at (25:3.55) {\small $s$};
  % angle
  \draw[->,seccolor,thick] (1.3,0) arc (0:50:1.3);
  \node[seccolor] at (25:1.0) {\small $\theta$};
  % center
  \fill (0,0) circle (1.6pt);
  \node[below left] at (0,0) {\small $O$};
\end{tikzpicture}
\end{LTR}
\end{center}
\figcap{خاکہ ۱: زاویائی مقام \m{\theta}، قوس \m{s} اور نصف قطر \m{r}۔ خطِ حوالہ مثبت \textenglish{$x$} محور سے زاویہ \m{\theta} بناتا ہے۔}

ریڈین دو لمبائیوں کا تناسب ہونے کے سبب ایک خالص عدد ہے اور بے اِبعاد ہے۔ چونکہ \m{r} نصف قطر والے دائرے کا محیط \m{2\pi r} ہوتا ہے، اِس لیے پورے دائرے میں \m{2\pi} ریڈین ہوتے ہیں:
\[ 1\ \text{rev} = 360^{\circ} = 2\pi\ \text{rad}, \qquad 1\ \text{rad} \approx 57.3^{\circ} \approx 0.159\ \text{rev}. \]

\heading{زاویائی انتقال}

اگر جسم اپنا زاویائی مقام \m{\theta_1} سے \m{\theta_2} تک بدلے تو \textbf{زاویائی انتقال} (\textenglish{angular displacement}):
\[ \Delta\theta = \theta_2 - \theta_1. \]
اصول: گھڑی کی مخالف سمت میں گردش \textbf{مثبت} اور گھڑی کی سمت میں \textbf{منفی} ہوتی ہے۔ (یاد رکھنے کا آسان جملہ: ”گھڑیاں منفی ہوتی ہیں“۔)

\heading{زاویائی رفتار}

\m{\Delta t} وقفے میں \textbf{اوسط زاویائی رفتار}:
\[ \omega_{\text{avg}} = \frac{\theta_2-\theta_1}{t_2-t_1} = \frac{\Delta\theta}{\Delta t}. \]
\textbf{لمحاتی زاویائی رفتار} (\m{\omega}) اِسی تناسب کی وہ حد ہے جب \m{\Delta t \to 0}:
\[ \omega = \frac{d\theta}{dt}. \]
اِس کی اکائی عموماً \m{\mathrm{rad/s}} ہوتی ہے۔ زاویائی رفتار کی مقدار (\textenglish{magnitude}) کو \textbf{زاویائی سرعت} کہتے ہیں۔

\heading{زاویائی اسراع}

اگر زاویائی رفتار \m{\omega_1} سے \m{\omega_2} تک بدلے تو \textbf{اوسط زاویائی اسراع}:
\[ \alpha_{\text{avg}} = \frac{\omega_2-\omega_1}{t_2-t_1} = \frac{\Delta\omega}{\Delta t}, \qquad \alpha = \frac{d\omega}{dt}. \]
اکائی عموماً \m{\mathrm{rad/s^2}} ہوتی ہے۔

\heading{کیا زاویائی مقداریں سمتیہ ہیں؟}

زاویائی رفتار اور زاویائی اسراع کو \textbf{سمتیہ} (\textenglish{vectors}) کے طور پر دکھایا جا سکتا ہے: اِن کا رخ محورِ گردش کے ساتھ ساتھ ہوتا ہے، جسے \textbf{دائیں ہاتھ کے اصول} سے طے کیا جاتا ہے — دائیں ہاتھ کی اُنگلیاں گردش کی سمت میں موڑیں تو کھڑا انگوٹھا سمتیہ کا رخ بتاتا ہے۔ تاہم \textbf{زاویائی انتقال} (جب تک بہت چھوٹا نہ ہو) سمتیہ نہیں ہوتا، کیونکہ دو بڑے زاویائی انتقالوں کو جمع کرنے کا نتیجہ اُن کی ترتیب پر منحصر ہوتا ہے، اور سمتیہ جمع میں ترتیب سے فرق نہیں پڑنا چاہیے۔

\begin{center}
\begin{LTR}
\begin{tikzpicture}[scale=1.0]
  % vertical axis with omega vector
  \draw[gray!50,thin] (0,-1.8)--(0,2.6);
  % disk (ellipse) to suggest spinning record
  \draw[thick,seccolor,fill=lightfill] (0,0) ellipse (1.9 and 0.6);
  \fill (0,0) circle (1.6pt);
  % rotation direction arrow (counterclockwise seen from top)
  \draw[->,accent,very thick] (1.5,0.25) arc (20:320:1.6 and 0.5);
  % omega vector up
  \draw[->,accent2,very thick] (0,0)--(0,2.4) node[right]{\small $\vec{\omega}$};
  \node[seccolor] at (0,-1.05) {\small rotation axis};
\end{tikzpicture}
\end{LTR}
\end{center}
\figcap{خاکہ ۲: زاویائی رفتار کا سمتیہ \m{\vec{\omega}} محورِ گردش کے ساتھ ہوتا ہے۔ دائیں ہاتھ کا اصول اِس کا رخ طے کرتا ہے۔ (یہاں گردش اوپر سے دیکھنے پر گھڑی کی مخالف سمت میں ہے، اِس لیے \m{\vec{\omega}} اوپر کی طرف ہے۔)}

% ===================== 10-2 =====================
\bigheading{۱۰-۲ \ \ مستقل زاویائی اسراع کے ساتھ گردش}

جیسے انتقالی حرکت میں مستقل اسراع ایک اہم خاص صورت ہے (مثلاً آزاد گرتا جسم)، ویسے ہی گردش میں \textbf{مستقل زاویائی اسراع} (\m{\alpha = } مستقل) اہم ہے۔ اِس کے لیے حرکیاتی مساوات انتقالی مساوات کے بالکل متوازی ہیں — صرف \m{x \to \theta}، \m{v \to \omega}، \m{a \to \alpha} رکھ دیں:
\begin{align*}
\omega &= \omega_0 + \alpha t,\\
\theta - \theta_0 &= \omega_0 t + \tfrac{1}{2}\alpha t^2,\\
\omega^2 &= \omega_0^2 + 2\alpha(\theta-\theta_0),\\
\theta - \theta_0 &= \tfrac{1}{2}(\omega_0+\omega)\,t,\\
\theta - \theta_0 &= \omega t - \tfrac{1}{2}\alpha t^2.
\end{align*}
بنیادی مساوات صرف پہلی دو ہیں؛ باقی تینوں اِنھی سے اخذ ہوتی ہیں۔ کسی سوال میں ایسی مساوات چنیں جس میں صرف مطلوبہ نامعلوم متغیر ہو۔

\heading{مثال ۱: مستقل زاویائی اسراع}

\textbf{سوال:} ایک پہیہ سکون سے شروع ہوتا ہے اور \m{\alpha = 2.0\ \mathrm{rad/s^2}} کے مستقل اسراع سے گھومتا ہے۔ \m{t = 5.0\ \mathrm{s}} پر اِس کی زاویائی رفتار اور طے شدہ زاویہ کیا ہوں گے؟

\textbf{حل:} چونکہ \m{\omega_0 = 0}:
\[ \omega = \omega_0 + \alpha t = (2.0)(5.0) = 10\ \mathrm{rad/s}, \]
\[ \theta - \theta_0 = \tfrac{1}{2}\alpha t^2 = \tfrac{1}{2}(2.0)(5.0)^2 = 25\ \mathrm{rad}. \]

% ===================== 10-3 =====================
\bigheading{۱۰-۳ \ \ خطی اور زاویائی متغیرات کا تعلق}

جب کوئی صلب جسم محور کے گرد گھومتا ہے تو اُس کا ہر نقطہ اپنے اپنے دائرے میں گھومتا ہے۔ سب کی زاویائی رفتار \m{\omega} یکساں ہوتی ہے، مگر محور سے جتنا دور نقطہ ہو، اُس کی \textbf{خطی رفتار} اُتنی زیادہ ہوتی ہے۔ محور سے عمودی فاصلہ \m{r} اِن دونوں کو جوڑتا ہے۔

\subheading{مقام:}
\[ s = \theta r \qquad (\theta\ \text{ریڈین میں}). \]

\subheading{رفتار:} \m{s = \theta r} کا وقت کے لحاظ سے تفریق (\m{r} مستقل):
\[ v = \omega r. \]

\subheading{اسراع:} خطی اسراع کے دو اجزا ہوتے ہیں۔ \textbf{مماسی جزو} (\textenglish{tangential}) رفتار کی \emph{مقدار} بدلتا ہے:
\[ a_t = \alpha r. \]
\textbf{نصف قطری/مرکز مائل جزو} (\textenglish{radial / centripetal}) رفتار کی \emph{سمت} بدلتا ہے اور ہمیشہ محور کی طرف ہوتا ہے:
\[ a_r = \frac{v^2}{r} = \omega^2 r. \]
یکساں دائروی حرکت میں دور (\textenglish{period}):
\[ T = \frac{2\pi r}{v} = \frac{2\pi}{\omega}. \]

\begin{center}
\begin{LTR}
\begin{tikzpicture}[scale=1.0]
  \draw[thin,gray!50] (0,0) circle (2.8);
  \fill (0,0) circle (1.6pt);
  \node[below left] at (0,0) {\small $O$};
  % point P at 40 deg
  \coordinate (P) at (40:2.8);
  \fill[accent] (P) circle (2pt);
  \node[accent] at (40:3.1) {\small $P$};
  \draw[gray!60] (0,0)--(P) node[midway,above,sloped]{\small $r$};
  % tangential acceleration (perp to radius)
  \draw[->,accent2,very thick] (P)--($(P)!1.3cm!90:(0,0)$) node[above right]{\small $a_t$};
  % radial acceleration (toward center)
  \draw[->,accent,very thick] (P)--($(P)!1.3cm!(0,0)$) node[below left]{\small $a_r$};
\end{tikzpicture}
\end{LTR}
\end{center}
\figcap{خاکہ ۳: گھومتے جسم کے نقطہ \m{P} پر خطی اسراع کے دو اجزا — مماسی \m{a_t = \alpha r} (دائرے کے مماس پر) اور نصف قطری \m{a_r = \omega^2 r} (مرکز کی طرف)۔}

\heading{مثال ۲: خطی اور زاویائی متغیرات}

\textbf{سوال:} ایک قرص پر محور سے \m{r = 0.30\ \mathrm{m}} کے فاصلے پر ایک نقطہ ہے۔ کسی لمحے \m{\omega = 8.0\ \mathrm{rad/s}} اور \m{\alpha = 4.0\ \mathrm{rad/s^2}} ہیں۔ اِس نقطے کی خطی رفتار، مماسی اسراع اور نصف قطری اسراع نکالیں۔

\textbf{حل:}
\[ v = \omega r = (8.0)(0.30) = 2.4\ \mathrm{m/s}, \]
\[ a_t = \alpha r = (4.0)(0.30) = 1.2\ \mathrm{m/s^2}, \]
\[ a_r = \omega^2 r = (8.0)^2(0.30) = 19.2\ \mathrm{m/s^2}. \]

% ===================== 10-4 / 10-5 =====================
\bigheading{۱۰-۴ \ \ گردش کی حرکی توانائی اور گردشی جمود}

گھومتے جسم کے ہر ذرّے کی حرکی توانائی جوڑ کر کل حرکی توانائی ملتی ہے۔ \m{v_i = \omega r_i} رکھنے پر:
\[ K = \sum \tfrac{1}{2} m_i v_i^2 = \tfrac{1}{2}\Big(\sum m_i r_i^2\Big)\omega^2. \]
قوسین میں موجود مقدار بتاتی ہے کہ جسم کی کمیت محور کے گرد کیسے تقسیم ہے۔ اِسے \textbf{گردشی جمود} یا \textbf{عزمِ جمود} (\textenglish{rotational inertia / moment of inertia}) \m{I} کہتے ہیں:
\[ I = \sum m_i r_i^2 \quad\text{(منقطع ذرّات)}, \qquad I = \int r^2\, dm \quad\text{(متصل جسم)}. \]
اِس طرح گردش کی حرکی توانائی:
\[ \boxed{\,K = \tfrac{1}{2} I \omega^2\,}. \]
یہ انتقالی \m{K = \tfrac{1}{2} M v^2} کا زاویائی ہم منصب ہے: جہاں کمیت \m{M} تھی وہاں \m{I} ہے، اور رفتار کے مربع کی جگہ زاویائی رفتار کا مربع۔ \m{I} کی \textenglish{SI} اکائی \m{\mathrm{kg\cdot m^2}} ہے۔

ایک ہی کمیت محور کے قریب ہو تو گھمانا آسان (\m{I} چھوٹا)، دور ہو تو مشکل (\m{I} بڑا)۔ اِسی لیے ایک لمبی سلاخ کو اُس کے طولی محور پر گھمانا آسان ہے، مگر مرکز سے گزرتے عمودی محور پر مشکل۔

\heading{عام شکلوں کا گردشی جمود}

ذیل کی جدول چند عام جسموں کے \m{I} دیتی ہے (\m{M} کمیت، \m{R} نصف قطر، \m{L} لمبائی):

\begin{center}
\renewcommand{\arraystretch}{1.5}
\begin{tabular}{>{\raggedleft\arraybackslash}p{8.5cm} >{\centering\arraybackslash}p{4.5cm}}
\toprule
\textbf{جسم اور محور} & \textbf{گردشی جمود \m{I}} \\
\midrule
چھلا، مرکزی محور کے گرد & \m{MR^2} \\
ٹھوس اسطوانہ/قرص، مرکزی محور کے گرد & \m{\tfrac{1}{2}MR^2} \\
پتلی سلاخ، مرکز سے گزرتے عمودی محور & \m{\tfrac{1}{12}ML^2} \\
پتلی سلاخ، سرے سے گزرتے عمودی محور & \m{\tfrac{1}{3}ML^2} \\
ٹھوس کرہ، کسی قطر کے گرد & \m{\tfrac{2}{5}MR^2} \\
پتلا کروی خول، کسی قطر کے گرد & \m{\tfrac{2}{3}MR^2} \\
مستطیل تختہ (اضلاع \m{a,b})، مرکز سے عمودی محور & \m{\tfrac{1}{12}M(a^2+b^2)} \\
\bottomrule
\end{tabular}
\\[4pt]
{\small جدول ۱۰-۱: چند عام شکلوں کے گردشی جمود}
\end{center}

\heading{متوازی محور کا مسئلہ}

اگر کسی جسم کا گردشی جمود اُس کے \textbf{مرکزِ کمیت} سے گزرتے محور کے گرد \m{I_{\text{com}}} معلوم ہو، تو اُس کے متوازی، \m{h} فاصلے پر واقع کسی محور کے گرد جمود:
\[ \boxed{\,I = I_{\text{com}} + M h^2\,} \qquad (\text{متوازی محور کا مسئلہ}). \]
یہاں \m{h} دونوں محوروں کے درمیان عمودی فاصلہ ہے۔

\begin{center}
\begin{LTR}
\begin{tikzpicture}[scale=1.0]
  % blob
  \draw[thick,seccolor,fill=lightfill] plot[smooth cycle,tension=0.9]
    coordinates {(-2,0.2)(-1.2,1.3)(0.4,1.5)(1.7,0.7)(1.3,-0.9)(-0.4,-1.3)(-1.8,-0.7)};
  % com axis
  \fill (0,0) circle (1.8pt);
  \node[below left] at (0,0) {\small com};
  \draw[accent2,thick,dashed] (0,-1.9)--(0,1.9);
  \node[accent2] at (0,2.1) {\small com axis};
  % parallel axis at P
  \fill (2.6,0) circle (1.8pt);
  \node[below] at (2.6,-0.05) {\small $P$};
  \draw[accent,thick,dashed] (2.6,-1.9)--(2.6,1.9);
  \node[accent] at (2.6,2.1) {\small axis through $P$};
  % distance h
  \draw[<->,thick] (0,-1.6)--(2.6,-1.6) node[midway,below]{\small $h$};
\end{tikzpicture}
\end{LTR}
\end{center}
\figcap{خاکہ ۴: متوازی محور کا مسئلہ — مرکزِ کمیت سے گزرتے محور اور \m{P} سے گزرتے متوازی محور کے درمیان فاصلہ \m{h}۔}

\heading{مثال ۳: گردشی جمود اور حرکی توانائی}

\textbf{سوال:} ایک ٹھوس قرص (\m{M = 2.0\ \mathrm{kg}}، \m{R = 0.10\ \mathrm{m}}) \m{\omega = 30\ \mathrm{rad/s}} پر گھوم رہی ہے۔ اِس کا گردشی جمود اور حرکی توانائی نکالیں۔

\textbf{حل:}
\[ I = \tfrac{1}{2}MR^2 = \tfrac{1}{2}(2.0)(0.10)^2 = 0.010\ \mathrm{kg\cdot m^2}, \]
\[ K = \tfrac{1}{2}I\omega^2 = \tfrac{1}{2}(0.010)(30)^2 = 4.5\ \mathrm{J}. \]

\heading{مثال ۴: متوازی محور کا اطلاق}

\textbf{سوال:} پتلی سلاخ کا مرکز سے گزرتے عمودی محور کے گرد جمود \m{\tfrac{1}{12}ML^2} ہے۔ سرے سے گزرتے متوازی محور کے گرد جمود نکالیں۔

\textbf{حل:} سرا مرکز سے \m{h = L/2} پر ہے، اِس لیے:
\[ I = I_{\text{com}} + Mh^2 = \tfrac{1}{12}ML^2 + M\left(\tfrac{L}{2}\right)^2 = \tfrac{1}{12}ML^2 + \tfrac{1}{4}ML^2 = \tfrac{1}{3}ML^2, \]
جو جدول کی قدر سے میل کھاتا ہے۔

% ===================== 10-6 =====================
\bigheading{۱۰-۶ \ \ مروڑ}

\textbf{مروڑ} (\textenglish{torque}) کسی جسم پر، کسی محور کے گرد، قوت کے ذریعے پیدا ہونے والا \emph{گھمانے} یا \emph{موڑنے} کا عمل ہے۔ فرض کریں قوت \m{\vec{F}} نقطہ \m{P} پر لگتی ہے، جس کا محور سے مقام سمتیہ \m{\vec{r}} ہے، اور \m{\vec{r}} و \m{\vec{F}} کے درمیان زاویہ \m{\phi} ہے۔ مروڑ کی مقدار:
\[ \tau = r F \sin\phi = r F_t = r_\perp F, \]
جہاں \m{F_t = F\sin\phi} قوت کا \m{\vec{r}} پر عمودی (مماسی) جزو ہے، اور \m{r_\perp = r\sin\phi} \textbf{بازوئے مروڑ} (\textenglish{moment arm}) — یعنی محور سے قوت کے \textbf{خطِ عمل} تک عمودی فاصلہ — ہے۔

\begin{center}
\begin{LTR}
\begin{tikzpicture}[scale=1.1]
  % axis at O
  \fill (0,0) circle (1.8pt);
  \node[below left] at (0,0) {\small $O$ (axis)};
  % position vector r to P
  \coordinate (P) at (35:3.2);
  \draw[->,seccolor,very thick] (0,0)--(P) node[midway,below right]{\small $\vec{r}$};
  \fill[accent] (P) circle (2pt);
  \node[above left] at (P) {\small $P$};
  % force F at P, at angle phi from r
  \coordinate (Fdir) at ($(P)+(110:2.4)$);
  \draw[->,accent,very thick] (P)--(Fdir) node[above]{\small $\vec{F}$};
  % line of action (extend force both ways, dashed)
  \draw[accent!60,dashed] ($(P)!-1.2cm!(Fdir)$)--($(P)!1.2cm!(Fdir)$);
  \node[accent!70] at ($(P)!-1.5cm!(Fdir)$) {\small line of action};
  % moment arm: perpendicular from O to line of action
  \coordinate (Q) at ($(P)!(0,0)!(Fdir)$);
  \draw[->,accent2,very thick] (0,0)--(Q) node[midway,above left]{\small $r_\perp$};
  \draw[accent2] ($(Q)+(0.0,0)$) rectangle ++(0.18,0.18);
  % angle phi between r and F at P
  \draw[seccolor,thick] (P)+(215:0.7) arc (215:110:0.7);
  \node[seccolor] at ($(P)+(160:1.05)$) {\small $\phi$};
\end{tikzpicture}
\end{LTR}
\end{center}
\figcap{خاکہ ۵: مروڑ \m{\tau = rF\sin\phi}۔ \m{\vec{r}} مقام سمتیہ، \m{\vec{F}} قوت، \m{\phi} اُن کے درمیان زاویہ، اور \m{r_\perp} بازوئے مروڑ (محور سے خطِ عمل تک عمودی فاصلہ)۔}

اکائی \textbf{نیوٹن-میٹر} (\m{\mathrm{N\cdot m}}) ہے۔ (احتیاط: کام کی اکائی بھی نیوٹن-میٹر ہے، مگر مروڑ اور کام بالکل مختلف مقداریں ہیں؛ کام جُول میں لکھا جاتا ہے، مروڑ کبھی نہیں۔) ”گھڑیاں منفی ہوتی ہیں“ یہاں بھی لاگو ہے: گھڑی کی مخالف سمت میں گھمانے والا مروڑ مثبت، گھڑی کی سمت میں منفی۔ کئی مروڑ ہوں تو خالص مروڑ \m{\tau_{\text{net}}} اُن کا (نشانی سمیت) مجموعہ ہوتا ہے۔

% ===================== 10-7 =====================
\bigheading{۱۰-۷ \ \ گردش کے لیے نیوٹن کا دوسرا قانون}

جیسے انتقالی حرکت میں \m{F_{\text{net}} = M a}، ویسے ہی گردش میں خالص مروڑ کو زاویائی اسراع سے جوڑا جاتا ہے۔ \m{F \to \tau}، \m{M \to I}، \m{a \to \alpha} رکھنے پر:
\[ \boxed{\,\tau_{\text{net}} = I\alpha\,} \qquad (\text{ریڈین میں}). \]

\textbf{مختصر ثبوت:} ایک بے کمیت سلاخ کے سرے پر کمیت \m{m} والا ذرّہ لیں جو \m{r} نصف قطر کے دائرے میں گھومتا ہے۔ مماسی قوت کے لیے \m{F_t = m a_t}، اور مروڑ \m{\tau = F_t\, r = m a_t r}۔ \m{a_t = \alpha r} رکھنے پر:
\[ \tau = m(\alpha r)r = (m r^2)\alpha = I\alpha. \]
کسی بھی صلب جسم کو ایسے ذرّات کا مجموعہ سمجھ کر یہی نتیجہ \m{\tau_{\text{net}} = I\alpha} حاصل ہوتا ہے۔

\heading{مثال ۵: مروڑ اور زاویائی اسراع}

\textbf{سوال:} کسی جسم پر خالص مروڑ \m{5.0\ \mathrm{N\cdot m}} لگتا ہے اور اُس کا گردشی جمود \m{2.5\ \mathrm{kg\cdot m^2}} ہے۔ زاویائی اسراع کیا ہوگا؟

\textbf{حل:}
\[ \alpha = \frac{\tau_{\text{net}}}{I} = \frac{5.0}{2.5} = 2.0\ \mathrm{rad/s^2}. \]

% ===================== 10-8 =====================
\bigheading{۱۰-۸ \ \ کام اور گردشی حرکی توانائی}

جب کوئی مروڑ کسی جسم کو محور کے گرد گھماتا ہے تو وہ کام کرتا ہے، جس سے جسم کی گردشی حرکی توانائی بدلتی ہے۔ اِن مقداروں کی مساوات انتقالی حرکت کے متوازی ہیں:
\[ W = \int_{\theta_i}^{\theta_f} \tau\, d\theta, \qquad P = \frac{dW}{dt} = \tau\,\omega. \]
جب \m{\tau} مستقل ہو:
\[ W = \tau\,(\theta_f - \theta_i). \]
\textbf{کام-حرکی توانائی کا مسئلہ} (گردشی صورت):
\[ \Delta K = K_f - K_i = \tfrac{1}{2} I \omega_f^2 - \tfrac{1}{2} I \omega_i^2 = W. \]

\heading{مثال ۶: مروڑ کا کیا گیا کام}

\textbf{سوال:} ایک مستقل مروڑ \m{3.0\ \mathrm{N\cdot m}} کسی جسم کو \m{4.0\ \mathrm{rad}} گھماتا ہے۔ کیا گیا کام کتنا ہے؟ اگر اِسی لمحے \m{\omega = 10\ \mathrm{rad/s}} ہو تو طاقت کیا ہے؟

\textbf{حل:}
\[ W = \tau\,\Delta\theta = (3.0)(4.0) = 12\ \mathrm{J}, \qquad P = \tau\,\omega = (3.0)(10) = 30\ \mathrm{W}. \]

% ===================== Summary =====================
\bigheading{خلاصہ}

\textbf{گردشی متغیرات:} \m{\theta = s/r}؛ \m{\Delta\theta = \theta_2-\theta_1}؛ \m{\omega = d\theta/dt}؛ \m{\alpha = d\omega/dt}۔ گھڑی کی مخالف سمت مثبت۔

\textbf{مستقل زاویائی اسراع:} انتقالی حرکیات کے متوازی پانچ مساوات، جن کی بنیاد \m{\omega = \omega_0 + \alpha t} اور \m{\theta-\theta_0 = \omega_0 t + \tfrac{1}{2}\alpha t^2} ہیں۔

\textbf{خطی–زاویائی تعلق:} \m{s = \theta r}، \m{v = \omega r}، \m{a_t = \alpha r}، \m{a_r = \omega^2 r}، \m{T = 2\pi/\omega}۔

\textbf{گردشی جمود اور توانائی:} \m{I = \sum m_i r_i^2 = \int r^2 dm}؛ \m{K = \tfrac{1}{2}I\omega^2}؛ متوازی محور: \m{I = I_{\text{com}} + Mh^2}۔

\textbf{مروڑ:} \m{\tau = rF\sin\phi = r_\perp F}؛ اکائی \m{\mathrm{N\cdot m}}۔

\textbf{گردش کا نیوٹن دوسرا قانون:} \m{\tau_{\text{net}} = I\alpha}۔

\textbf{کام اور طاقت:} \m{W = \int \tau\, d\theta}؛ مستقل مروڑ پر \m{W = \tau\,\Delta\theta}؛ \m{P = \tau\,\omega}؛ \m{\Delta K = W}۔

\end{document}
